\documentclass[12pt]{beamer}

\usetheme{Boadilla}
\usecolortheme{default}

\usepackage[utf8]{inputenc}
\usepackage[T1]{fontenc}
\usepackage[brazil]{babel}
\usepackage{amsmath, amssymb}
\usepackage{booktabs}

% Informações
\title{O que causa variação na abundância?}
\subtitle{Construindo o primeiro modelo populacional}
\author{Modelagem Matemática}
\institute{Matemática Aplicada}
\date{Aula 02}

\begin{document}

% ── CAPA ──────────────────────────────────────────
\begin{frame}
  \titlepage
\end{frame}

% ══════════════════════════════════════════════════
\section{Ponte com a aula anterior}
% ══════════════════════════════════════════════════

\begin{frame}{Onde estamos}

  Na aula anterior chegamos a duas conclusões:

  \bigskip
  \begin{block}{Conclusão 1 — A pergunta de Krebs}
    Quais interações determinam \emph{onde} os organismos
    ocorrem e em \emph{que quantidade}?
  \end{block}

  \begin{exampleblock}{Conclusão 2 — A linguagem necessária}
    Para responder a essa pergunta, precisamos de uma linguagem
    que represente como \textbf{quantidades mudam em função de interações}.
    Essa linguagem é o \textbf{cálculo diferencial}.
  \end{exampleblock}

  \bigskip
  Sabemos que queremos explicar $N(t)$ --- a abundância ao longo do tempo.

  \bigskip
  \textbf{Pergunta desta aula:} o que faz $N$ variar?

\end{frame}

% ══════════════════════════════════════════════════
\section{Bloco 1 — Contabilidade populacional}
% ══════════════════════════════════════════════════

\begin{frame}{O que pode mudar o número de indivíduos?}

  Antes de qualquer equação --- uma pergunta simples:

  \bigskip
  \begin{alertblock}{Pergunta}
    Quais são todos os processos que podem \textbf{aumentar}
    ou \textbf{diminuir} o número de indivíduos em uma população?
  \end{alertblock}

  \bigskip
  Pense em uma população de capivaras no Pantanal.
  O que pode fazer esse número subir? O que pode fazê-lo cair?

\end{frame}

% ──────────────────────────────────────────────────
\begin{frame}{Quatro e apenas quatro processos}

  Não importa a espécie, o ambiente ou a escala de tempo.
  A variação em $N$ tem exatamente \textbf{quatro causas}:

  \bigskip
  \begin{columns}[t]
    \column{0.48\textwidth}
      \begin{exampleblock}{Aumentam $N$}
        \begin{itemize}
          \item \textbf{Nascimentos} ($B$)
          \item \textbf{Imigração} ($I$)
        \end{itemize}
      \end{exampleblock}
    \column{0.48\textwidth}
      \begin{alertblock}{Diminuem $N$}
        \begin{itemize}
          \item \textbf{Mortes} ($D$)
          \item \textbf{Emigração} ($E$)
        \end{itemize}
      \end{alertblock}
  \end{columns}

\end{frame}

% ──────────────────────────────────────────────────
\begin{frame}{Uma identidade contábil}

  Isso não é um modelo --- é uma \textbf{identidade}.
  Vale sempre, para qualquer população, sem nenhuma hipótese:

  \[
    \Delta N = B - D + I - E
  \]

  \bigskip
  \begin{block}{Observação importante}
    Uma identidade contábil \emph{descreve} o que acontece,
    mas não \emph{explica} por que acontece.
    Para explicar, precisamos entender o que determina
    $B$, $D$, $I$ e $E$.
  \end{block}

  \bigskip
  Essa será a tarefa das próximas aulas.
  Por enquanto, vamos simplificar.

\end{frame}

% ══════════════════════════════════════════════════
\section{Bloco 2 — População fechada}
% ══════════════════════════════════════════════════

\begin{frame}{Quando podemos ignorar migração?}

  \begin{alertblock}{Pergunta}
    Em quais situações podemos considerar que
    imigração e emigração são desprezíveis?
  \end{alertblock}

  \bigskip
  Exemplos:
  \begin{itemize}
    \item Uma população em uma ilha isolada
    \item Uma espécie endêmica de uma bacia hidrográfica
    \item Um experimento em laboratório
  \end{itemize}

\end{frame}

% ──────────────────────────────────────────────────
\begin{frame}{Hipótese 1 — População fechada}

  \begin{exampleblock}{Hipótese explícita}
    Assumimos que não há fluxo de indivíduos com outras populações:
    $I = 0$ e $E = 0$.
  \end{exampleblock}

  \bigskip
  A identidade contábil se reduz a:

  \[
    \Delta N = B - D
  \]

  \bigskip
  \begin{block}{Atenção}
    Toda hipótese simplificadora tem um custo.
    Ao assumir população fechada, perdemos a capacidade
    de modelar dispersão, recolonização e fluxo gênico.
    Ganhamos em tratabilidade.
  \end{block}

\end{frame}

% ──────────────────────────────────────────────────
\begin{frame}{Nova pergunta}

  Temos $\Delta N = B - D$.

  \bigskip
  \begin{alertblock}{Pergunta}
    $B$ e $D$ dependem de quê?
  \end{alertblock}

  \bigskip
  Uma população de 1000 capivaras produz mais filhotes
  que uma de 10 capivaras --- isso é óbvio.

  \bigskip
  Mas será que \emph{cada capivara} contribui da mesma forma
  nos dois casos?

  \bigskip
  O que seria mais informativo comparar?

\end{frame}

% ══════════════════════════════════════════════════
\section{Bloco 3 — Taxas per capita}
% ══════════════════════════════════════════════════

\begin{frame}{Taxas per capita}

  O que é constante --- se é que algo é constante --- não é
  o número total de nascimentos, mas a \textbf{contribuição por indivíduo}.

  \bigskip
  Definimos as \textbf{taxas per capita}:

  \[
    b = \frac{B}{N} \qquad d = \frac{D}{N}
  \]

  onde $b$ é a taxa de natalidade per capita e
  $d$ é a taxa de mortalidade per capita.

  \bigskip
  \begin{block}{Intuição}
    $b = 0{,}3$ significa que, em média, cada indivíduo
    contribui com 0,3 nascimentos por unidade de tempo.
  \end{block}

\end{frame}

% ──────────────────────────────────────────────────
\begin{frame}{Hipótese 2 — Taxas per capita constantes}

  \begin{exampleblock}{Hipótese explícita}
    Assumimos que $b$ e $d$ são constantes ---
    não dependem de $N$, do tempo, nem do ambiente.
  \end{exampleblock}

  \bigskip
  Substituindo $B = bN$ e $D = dN$ em $\Delta N = B - D$:

  \[
    \frac{dN}{dt} = bN - dN = (b - d)\,N
  \]

  \bigskip
  Definindo $r = b - d$ como a \textbf{taxa intrínseca de crescimento}:

  \[
    \boxed{\frac{dN}{dt} = rN}
  \]

\end{frame}

% ──────────────────────────────────────────────────
\begin{frame}{O que $r$ significa?}

  \begin{itemize}
    \item $r > 0$ \quad natalidade supera mortalidade
      --- população \textbf{cresce}
    \bigskip
    \item $r = 0$ \quad natalidade igual à mortalidade
      --- população \textbf{estável}
    \bigskip
    \item $r < 0$ \quad mortalidade supera natalidade
      --- população \textbf{declina}
  \end{itemize}

  \bigskip
  \begin{alertblock}{Pergunta}
    Uma população pode crescer indefinidamente?
    Essa hipótese é razoável na natureza?
  \end{alertblock}

\end{frame}

% ══════════════════════════════════════════════════
\section{Bloco 4 — Quebrando a hipótese}
% ══════════════════════════════════════════════════

\begin{frame}{Quando a hipótese falha?}

  A hipótese de taxas per capita constantes implica
  crescimento ilimitado. Mas observamos na natureza:

  \bigskip
  \begin{itemize}
    \item Populações que crescem rápido quando pequenas
      e desaceleram quando grandes
    \item Populações que se estabilizam em torno de
      um tamanho típico para o ambiente
    \item Populações que colapsam quando superam
      a capacidade do ambiente
  \end{itemize}

  \bigskip
  \begin{alertblock}{Pergunta}
    O que biologicamente explica essa desaceleração?
  \end{alertblock}

\end{frame}

% ──────────────────────────────────────────────────
\begin{frame}{Recursos são finitos}

  À medida que $N$ cresce:

  \bigskip
  \begin{itemize}
    \item Alimento por indivíduo \textbf{diminui} $\rightarrow$
      $b$ cai, $d$ sobe
    \item Espaço por indivíduo \textbf{diminui} $\rightarrow$
      competição aumenta
    \item Predação e doenças se tornam \textbf{mais frequentes}
      $\rightarrow$ mortalidade sobe
  \end{itemize}

  \bigskip
  \begin{exampleblock}{Conclusão}
    As taxas per capita \textbf{não são constantes} ---
    elas dependem de $N$.
    A hipótese 2 precisa ser revisada.
  \end{exampleblock}

\end{frame}

% ──────────────────────────────────────────────────
\begin{frame}{Hipótese 3 — Dependência de densidade}

  \begin{exampleblock}{Hipótese explícita}
    As taxas per capita diminuem linearmente com $N$.
    Existe um tamanho $K$ --- a \textbf{capacidade de suporte} ---
    no qual o crescimento líquido se anula.
  \end{exampleblock}

  \bigskip
  O fator $(1 - N/K)$ captura essa ideia:

  \bigskip
  \begin{itemize}
    \item Quando $N \ll K$: \quad $(1 - N/K) \approx 1$
      --- cresce quase como exponencial
    \item Quando $N = K/2$: \quad $(1 - N/K) = 1/2$
      --- crescimento reduzido à metade
    \item Quando $N = K$: \quad $(1 - N/K) = 0$
      --- crescimento nulo
    \item Quando $N > K$: \quad $(1 - N/K) < 0$
      --- população \textbf{declina}
  \end{itemize}

\end{frame}

% ──────────────────────────────────────────────────
\begin{frame}{O segundo modelo}

  Incorporando a dependência de densidade:

  \[
    \boxed{\frac{dN}{dt} = rN\!\left(1 - \frac{N}{K}\right)}
  \]

  \bigskip
  Cada símbolo é uma \textbf{hipótese explícita}:

  \bigskip
  \begin{itemize}
    \item $N$ --- abundância (o que queremos explicar)
    \item $r$ --- taxa intrínseca de crescimento (sem limitação)
    \item $K$ --- capacidade de suporte (limite imposto pelo ambiente)
    \item $(1 - N/K)$ --- fração de recursos ainda disponíveis
  \end{itemize}

\end{frame}

% ══════════════════════════════════════════════════
\section{Síntese}
% ══════════════════════════════════════════════════

\begin{frame}{Síntese — Como chegamos até aqui}

  \begin{enumerate}
    \item \textbf{Identidade contábil} \hfill $\Delta N = B - D + I - E$
    \bigskip
    \item \textbf{Hipótese 1: população fechada} \hfill $\Delta N = B - D$
    \bigskip
    \item \textbf{Hipótese 2: taxas per capita constantes}
      \hfill $\dfrac{dN}{dt} = rN$
    \bigskip
    \item \textbf{Hipótese 2 falha} \hfill recursos são finitos
    \bigskip
    \item \textbf{Hipótese 3: dependência de densidade}
      \hfill $\dfrac{dN}{dt} = rN\!\left(1 - \dfrac{N}{K}\right)$
  \end{enumerate}

\end{frame}

% ──────────────────────────────────────────────────
\begin{frame}{O modelo é um argumento}

  \begin{block}{Para refletir}
    Cada equação que escrevemos é a consequência de
    hipóteses explícitas sobre mecanismos biológicos.
    Mudar uma hipótese muda a equação --- e muda a predição.
  \end{block}

  \bigskip
  \begin{exampleblock}{Próxima aula}
    Com os dois modelos em mãos, vamos analisar seu comportamento:
    \begin{itemize}
      \item Quais são os equilíbrios de cada modelo?
      \item Eles são estáveis ou instáveis?
      \item O que os dados reais dizem sobre $r$ e $K$?
    \end{itemize}
  \end{exampleblock}

\end{frame}

\end{document}
